\documentclass[12pt]{article}
\usepackage[T1]{fontenc}
\usepackage[utf8]{inputenc}
\usepackage{lmodern}
\usepackage{textcomp}
\usepackage{enumitem}
\usepackage{geometry}
\usepackage{graphicx}
\usepackage{amsmath, amssymb}
\geometry{margin=1in}
\begin{document}
\begin{center}
History - Undergraduate Level Assessment 4 \
Total Marks: 40 \
Answer all questions.
\end{center}

\section*{Multiple Choice (10 marks)}
\begin{enumerate}[label=\arabic*.]
\item Which of the following crops was NOT an important component of Hopewell subsistence?
\begin{enumerate}[label=(\alph*)]
    \item squash
    \item knotweed
    \item beans
    \item sunflower
    \item maize
\end{enumerate}
\item Australian and America megafauna were probably wiped out by:
\begin{enumerate}[label=(\alph*)]
    \item natural disasters such as volcanic eruptions or earthquakes.
    \item both humans and environmental changes.
    \item the spread of invasive plant species which altered their habitat.
    \item a sudden drastic climate change due to global warming.
    \item humans who carried diseases over the land bridge.
\end{enumerate}
\item Mayan art and iconography prominently feature depictions of:
\begin{enumerate}[label=(\alph*)]
    \item mutilation and violence.
    \item stars and comets.
    \item infants and children.
    \item modern technologies.
    \item domesticated animals.
\end{enumerate}
\item How did Aurignacian technology differ from Mousterian technology?
\begin{enumerate}[label=(\alph*)]
    \item Mousterian technology was focused on the production of scrapers.
    \item Aurignacian technology was primarily used for creating stone tools.
    \item Aurignacian technology was less advanced than Mousterian technology.
    \item Mousterian technology produced more usable blade surface.
    \item Aurignacian technology produced more usable blade surface.
\end{enumerate}
\item There is evidence that premodern Homo sapiens existed:
\begin{enumerate}[label=(\alph*)]
    \item 500,000 to 100,000 years ago.
    \item 800,000 to 200,000 years ago.
    \item 200,000 to 50,000 years ago.
    \item 700,000 to 40,000 years ago.
    \item 600,000 to 30,000 years ago.
\end{enumerate}
\end{enumerate}

\section*{True / False (10 marks)}
\textit{Answer True or False}
\begin{enumerate}[label=\arabic*.]
\setcounter{enumi}{5}
\item True or False: Dell assembled computers for the EMEA market at the Limerick facility in the Republic of Ireland, and once employed about 4,500 people in that country.
\item True or False: Edward Weston (1850-1936), a British-born electrical engineer, industrialist and founder of the US-based Weston Electrical Instrument Corporation, with the Weston model 617, one of the earliest photo-electric exposure meters, in August 1932.
\item True or False: The Jawa Dwipa Hindu kingdom in Java and Sumatra existed around 200 BCE. The history of the Malay-speaking world began with the advent of Indian influence, which dates back to at least the 3 rd century BCE. Indian traders came to the archipelago both for its abundant forest and maritime products and to trade with merchants from China, who also discovered the Malay world at an early date.
\item True or False: However, the army was hopeless in battle against the western forces, particularly against the young Mahmud of Ghazni.
\item True or False: Millions of people moved to and from British colonies, with large numbers of Indians emigrating to other parts of the empire, such as Malaysia and Fiji, and Chinese people to Malaysia, Singapore and the Caribbean. The demographics of Britain itself was changed after the Second World War owing to immigration to Britain from its former colonies.
\end{enumerate}

\section*{Long Form Response (20 marks)}
\textit{Answer each question with a clear and complete explanation.}
\begin{enumerate}[label=\arabic*.]
\setcounter{enumi}{10}
\item Read the following passage and answer the question: Ancient Greece had traditionally been a fractious collection of fiercely independent city-states. After the Peloponnesian War (431-404 BC), Greece had fallen under a Spartan hegemony, in which Sparta was pre-eminent but not all-powerful. Spartan hegemony was succeeded by a Theban one after the Battle of Leuctra (371 BC), but after the Battle of Mantinea (362 BC), all of Greece was so weakened that no one state could claim pre-eminence. It was against this backdrop, that the ascendancy of Macedon began, under king Philip II. Macedon was located at the periphery of the Greek world, and although its royal family claimed Greek descent, the Macedonians themselves were looked down upon as semi-barbaric by the rest of the Greeks. However, Macedon had a relatively strong and centralised government... Question: What hegemoney replaced Sparta after the Battle of Leuctra?
\item Read the following passage and answer the question: In 2000, Madonna starred in the film The Next Best Thing, and contributed two songs to the film's soundtrack; "Time Stood Still" and a cover of Don McLean's 1971 song "American Pie". She released her eighth studio album, Music, in September 2000. It featured elements from the electronica-inspired Ray of Light era, and like its predecessor, received acclaim from critics. Collaborating with French producer Mirwais Ahmadzaï, Madonna commented: "I love to work with the weirdos that no one knows about-the people who have raw talent and who are making music unlike anyone else out there. Music is the future of sound." Stephen Thomas Erlewine from AllMusic felt that "Music blows by in a kaleidoscopic rush of color, technique, style and substance. It has so many depth and layers that it's easily as... Question: What was Madonna's eighth album called?
\end{enumerate}

\vfill
\noindent\textit{End of Paper}
\end{document}